% !PRINTED 419
\weberpage{419}
%<<CONT
besonderen Werte des Arguments $\omega$, die wir im vorigen Paragraphen betrachtet haben. Von diesen hängen die Moduln der elliptischen Funktionen ab, die eine komplexe Multiplikation zulassen. Sie heißen nach \gesperrt{Kronecker} \gesperrt{singuläre Moduln}. Wir werden dementsprechend auch von den \gesperrt{singulären Werten der Modulfunktionen}, insbesondere von den \gesperrt{singulären Invarianten} $j(\omega)$ sprechen, und verstehen darunter die Werte, die diese Funktionen annehmen, wenn $\omega$ die \gesperrt{Wurzel} (mit positiv imaginärem Bestandteil) \gesperrt{einer quadratischen Form mit negativer Diskriminante} ist.

Wenn $\omega$ der Gleichung (3), § 114, genügt und $a, b, c, \partial, n$ wie im vorigen Paragraphen bestimmt sind, so folgt zunächst:

\weq{1}{j(\omega) = j\left(\frac{c + \partial\,\omega}{a + b\,\omega}\right),}

und wenn also

\weq{2}{F_n(u, v) = 0}

die zum Transformationsgrad $n$ gehörige Invariantengleichung ist (§ 69), so ist (2) befriedigt, wenn $a, b, c, \partial$ ohne gemeinsamen Teiler sind und

\weq{3}{u = j(\omega), \quad v = j\left(\frac{c + \partial\,\omega}{a + b\,\omega}\right) = u}

gesetzt wird. Danach gelangen wir zu dem ersten Hauptsatz dieser Theorie:

III. \gesperrt{Der singuläre Wert}
\[
u = j(\omega)
\]
\gesperrt{ist eine Wurzel der algebraischen Gleichung:}

\weq{4}{F_n(u, u) = 0.}

Ist umgekehrt $u$ eine Wurzel der Gleichung (4), so ist, wenn $\omega$ aus der Gleichung $j(\omega) = u$ bestimmt wird, einer von den der Gleichung (2) genügenden Werten von $v$ gleichfalls $= u$, und es besteht also eine Gleichung von der Form (1); diese hat nach § 53 zur Folge, daß $\omega$ mit $\dfrac{c + \partial\,\omega}{a + b\,\omega}$ äquivalent ist, und daraus ergibt sich eine Gleichung von der Form § 114, (3). Also haben wir die Ergänzung zu III.:

IV. \gesperrt{Jede Wurzel von (4) ist eine singuläre Invariante.}
